%% Generated by Sphinx.
\def\sphinxdocclass{jupyterBook}
\documentclass[letterpaper,10pt,english]{jupyterBook}
\ifdefined\pdfpxdimen
   \let\sphinxpxdimen\pdfpxdimen\else\newdimen\sphinxpxdimen
\fi \sphinxpxdimen=.75bp\relax
\ifdefined\pdfimageresolution
    \pdfimageresolution= \numexpr \dimexpr1in\relax/\sphinxpxdimen\relax
\fi
%% let collapsible pdf bookmarks panel have high depth per default
\PassOptionsToPackage{bookmarksdepth=5}{hyperref}
%% turn off hyperref patch of \index as sphinx.xdy xindy module takes care of
%% suitable \hyperpage mark-up, working around hyperref-xindy incompatibility
\PassOptionsToPackage{hyperindex=false}{hyperref}
%% memoir class requires extra handling
\makeatletter\@ifclassloaded{memoir}
{\ifdefined\memhyperindexfalse\memhyperindexfalse\fi}{}\makeatother

\PassOptionsToPackage{booktabs}{sphinx}
\PassOptionsToPackage{colorrows}{sphinx}

\PassOptionsToPackage{warn}{textcomp}

\catcode`^^^^00a0\active\protected\def^^^^00a0{\leavevmode\nobreak\ }
\usepackage{cmap}
\usepackage{fontspec}
\defaultfontfeatures[\rmfamily,\sffamily,\ttfamily]{}
\usepackage{amsmath,amssymb,amstext}
\usepackage{polyglossia}
\setmainlanguage{english}



\setmainfont{FreeSerif}[
  Extension      = .otf,
  UprightFont    = *,
  ItalicFont     = *Italic,
  BoldFont       = *Bold,
  BoldItalicFont = *BoldItalic
]
\setsansfont{FreeSans}[
  Extension      = .otf,
  UprightFont    = *,
  ItalicFont     = *Oblique,
  BoldFont       = *Bold,
  BoldItalicFont = *BoldOblique,
]
\setmonofont{FreeMono}[
  Extension      = .otf,
  UprightFont    = *,
  ItalicFont     = *Oblique,
  BoldFont       = *Bold,
  BoldItalicFont = *BoldOblique,
]



\usepackage[Bjarne]{fncychap}
\usepackage[,numfigreset=1,mathnumfig]{sphinx}

\fvset{fontsize=\small}
\usepackage{geometry}


% Include hyperref last.
\usepackage{hyperref}
% Fix anchor placement for figures with captions.
\usepackage{hypcap}% it must be loaded after hyperref.
% Set up styles of URL: it should be placed after hyperref.
\urlstyle{same}


\usepackage{sphinxmessages}



        % Start of preamble defined in sphinx-jupyterbook-latex %
         \usepackage[Latin,Greek]{ucharclasses}
        \usepackage{unicode-math}
        % fixing title of the toc
        \addto\captionsenglish{\renewcommand{\contentsname}{Contents}}
        \hypersetup{
            pdfencoding=auto,
            psdextra
        }
        % End of preamble defined in sphinx-jupyterbook-latex %
        

\title{Enquête Science Ouverte}
\date{Sep 25, 2026}
\release{}
\author{BU Paris 8}
\newcommand{\sphinxlogo}{\vbox{}}
\renewcommand{\releasename}{}
\makeindex
\begin{document}

\pagestyle{empty}
\sphinxmaketitle
\pagestyle{plain}
\sphinxtableofcontents
\pagestyle{normal}
\phantomsection\label{\detokenize{intro::doc}}


\sphinxAtStartPar
Dans un contexte de transformation profonde des pratiques de recherche sous l’effet du numérique et du mouvement international en faveur de la science ouverte, l’université engage une
démarche structurée de réflexion et d’action autour de la gestion, du partage et de la
valorisation des productions scientifiques. Cette dynamique s’inscrit dans la volonté de construire,
à moyen et long terme, une politique d’établissement cohérente en matière de science ouverte,
fondée sur une compréhension fine des besoins de sa communauté et sur la mise en œuvre
progressive d’une offre de services adaptée.

\sphinxAtStartPar
C’est dans ce cadre qu’une enquête est conduite auprès des chercheurs, des personnels de laboratoire et des doctorants, afin de mieux connaître leurs pratiques numériques, leurs environnements de travail et leurs besoins d’accompagnement et de formation autour de plusieurs thématiques clés :
\begin{itemize}
\item {} 
\sphinxAtStartPar
la science ouverte et ses principes ;

\item {} 
\sphinxAtStartPar
le cycle de vie des données de la recherche (collecte, traitement, gestion, documentation, ouverture et partage) ;

\item {} 
\sphinxAtStartPar
les pratiques de publication en accès ouvert et la diffusion des résultats scientifiques.

\end{itemize}

\sphinxstepscope


\chapter{Connaissance et pratique de la Science Ouverte}
\label{\detokenize{script_analyse/stat_descriptive/connaissance_science_ouverte:connaissance-et-pratique-de-la-science-ouverte}}\label{\detokenize{script_analyse/stat_descriptive/connaissance_science_ouverte::doc}}
\sphinxAtStartPar
Les premières questions interrogeaient les participants sur leurs connaissances et leurs pratiques de la Science Ouverte, thème central de l’enquête. Le questionnaire commencait ainsi par demander aux participants de préciser leur degré de familiarité avec les principes de la Science Ouverte (\hyperref[\detokenize{script_analyse/stat_descriptive/connaissance_science_ouverte:principe-so}]{Table \ref{\detokenize{script_analyse/stat_descriptive/connaissance_science_ouverte:principe-so}}}\}). 56\% ont répond “oui, un peu”, 32 “oui, tout à fait” et 12\% “non”. Les pratiques étaient ensuite abordaient à travers des questions sur le dépôt des publications sur Hal et Octaviana (\hyperref[\detokenize{script_analyse/stat_descriptive/connaissance_science_ouverte:depot-hal}]{Table \ref{\detokenize{script_analyse/stat_descriptive/connaissance_science_ouverte:depot-hal}}} et \hyperref[\detokenize{script_analyse/stat_descriptive/connaissance_science_ouverte:depot-octaviana}]{Table \ref{\detokenize{script_analyse/stat_descriptive/connaissance_science_ouverte:depot-octaviana}}}), et l’ouverture de données de recherche (\hyperref[\detokenize{script_analyse/stat_descriptive/connaissance_science_ouverte:opened-data}]{Table \ref{\detokenize{script_analyse/stat_descriptive/connaissance_science_ouverte:opened-data}}}).

\sphinxAtStartPar
40 \% des personnes interrogées ont déposé plusieurs fois leurs publications sur Hal et 35\% le font systématiquement. 13\%  ne l’ont jamais fais. Concernant Octaviana, un peu moins de 10\% des répondants ont déjà déposé des productions sur la bibliothèque numérique. Ces productions sont des captations d’événements pour trois d’entre elles et des travaux universitaires et des publications pour les huit autres. En tout une personne interrogée sur cinq n’a jamais entendu parlé d’Octaviana.


\begin{savenotes}\sphinxattablestart
\sphinxthistablewithglobalstyle
\centering
\sphinxcapstartof{table}
\sphinxthecaptionisattop
\sphinxcaption{Connaissance des principes de la Science Ouverte}\label{\detokenize{script_analyse/stat_descriptive/connaissance_science_ouverte:principe-so}}
\sphinxaftertopcaption
\begin{tabulary}{\linewidth}[t]{TTT}
\sphinxtoprule
\sphinxstyletheadfamily 
\sphinxAtStartPar
Principes de la science ouverte
&\sphinxstyletheadfamily 
\sphinxAtStartPar
nb
&\sphinxstyletheadfamily 
\sphinxAtStartPar
freq
\\
\sphinxmidrule
\sphinxtableatstartofbodyhook
\sphinxAtStartPar
Non
&
\sphinxAtStartPar
14
&
\sphinxAtStartPar
11,8
\\
\sphinxhline
\sphinxAtStartPar
Oui, tout à fait
&
\sphinxAtStartPar
38
&
\sphinxAtStartPar
31,9
\\
\sphinxhline
\sphinxAtStartPar
Oui, un peu
&
\sphinxAtStartPar
67
&
\sphinxAtStartPar
56,3
\\
\sphinxhline
\sphinxAtStartPar
Total
&
\sphinxAtStartPar
119
&
\sphinxAtStartPar
100,0
\\
\sphinxbottomrule
\end{tabulary}
\sphinxtableafterendhook\par
\sphinxattableend\end{savenotes}


\begin{savenotes}\sphinxattablestart
\sphinxthistablewithglobalstyle
\centering
\sphinxcapstartof{table}
\sphinxthecaptionisattop
\sphinxcaption{Le dépôt dans HAL}\label{\detokenize{script_analyse/stat_descriptive/connaissance_science_ouverte:depot-hal}}
\sphinxaftertopcaption
\begin{tabulary}{\linewidth}[t]{TTT}
\sphinxtoprule
\sphinxstyletheadfamily 
\sphinxAtStartPar
Dépôt dans Hal
&\sphinxstyletheadfamily 
\sphinxAtStartPar
nb
&\sphinxstyletheadfamily 
\sphinxAtStartPar
freq
\\
\sphinxmidrule
\sphinxtableatstartofbodyhook
\sphinxAtStartPar
Non
&
\sphinxAtStartPar
15
&
\sphinxAtStartPar
12,6
\\
\sphinxhline
\sphinxAtStartPar
Oui, plusieurs fois
&
\sphinxAtStartPar
48
&
\sphinxAtStartPar
40,3
\\
\sphinxhline
\sphinxAtStartPar
Oui, systématiquement
&
\sphinxAtStartPar
41
&
\sphinxAtStartPar
34,5
\\
\sphinxhline
\sphinxAtStartPar
Oui, une fois
&
\sphinxAtStartPar
15
&
\sphinxAtStartPar
12,6
\\
\sphinxhline
\sphinxAtStartPar
Total
&
\sphinxAtStartPar
119
&
\sphinxAtStartPar
100,0
\\
\sphinxbottomrule
\end{tabulary}
\sphinxtableafterendhook\par
\sphinxattableend\end{savenotes}


\begin{savenotes}\sphinxattablestart
\sphinxthistablewithglobalstyle
\centering
\sphinxcapstartof{table}
\sphinxthecaptionisattop
\sphinxcaption{Le dépôt dans Octaviana}\label{\detokenize{script_analyse/stat_descriptive/connaissance_science_ouverte:depot-octaviana}}
\sphinxaftertopcaption
\begin{tabulary}{\linewidth}[t]{TTT}
\sphinxtoprule
\sphinxstyletheadfamily 
\sphinxAtStartPar
Dépot sur Octaviana
&\sphinxstyletheadfamily 
\sphinxAtStartPar
nb
&\sphinxstyletheadfamily 
\sphinxAtStartPar
freq
\\
\sphinxmidrule
\sphinxtableatstartofbodyhook
\sphinxAtStartPar
Non
&
\sphinxAtStartPar
108
&
\sphinxAtStartPar
90,8
\\
\sphinxhline
\sphinxAtStartPar
Oui
&
\sphinxAtStartPar
11
&
\sphinxAtStartPar
9,2
\\
\sphinxhline
\sphinxAtStartPar
Total
&
\sphinxAtStartPar
119
&
\sphinxAtStartPar
100,0
\\
\sphinxbottomrule
\end{tabulary}
\sphinxtableafterendhook\par
\sphinxattableend\end{savenotes}


\begin{savenotes}\sphinxattablestart
\sphinxthistablewithglobalstyle
\centering
\sphinxcapstartof{table}
\sphinxthecaptionisattop
\sphinxcaption{Ouverture des données}\label{\detokenize{script_analyse/stat_descriptive/connaissance_science_ouverte:opened-data}}
\sphinxaftertopcaption
\begin{tabulary}{\linewidth}[t]{TTTT}
\sphinxtoprule
\sphinxstyletheadfamily 
\sphinxAtStartPar
Ouverture des données
&\sphinxstyletheadfamily 
\sphinxAtStartPar
nb
&\sphinxstyletheadfamily 
\sphinxAtStartPar
total
&\sphinxstyletheadfamily 
\sphinxAtStartPar
freq
\\
\sphinxmidrule
\sphinxtableatstartofbodyhook
\sphinxAtStartPar
Non
&
\sphinxAtStartPar
86
&
\sphinxAtStartPar
119
&
\sphinxAtStartPar
72,3
\\
\sphinxhline
\sphinxAtStartPar
Autre, précisez
&
\sphinxAtStartPar
24
&
\sphinxAtStartPar
119
&
\sphinxAtStartPar
20,7
\\
\sphinxhline
\sphinxAtStartPar
Oui, sur Nakala
&
\sphinxAtStartPar
5
&
\sphinxAtStartPar
119
&
\sphinxAtStartPar
4,2
\\
\sphinxhline
\sphinxAtStartPar
Oui, sur Zenodo
&
\sphinxAtStartPar
4
&
\sphinxAtStartPar
119
&
\sphinxAtStartPar
3,4
\\
\sphinxhline
\sphinxAtStartPar
Oui, sur Data Paris 8
&
\sphinxAtStartPar
1
&
\sphinxAtStartPar
119
&
\sphinxAtStartPar
0,8
\\
\sphinxbottomrule
\end{tabulary}
\sphinxtableafterendhook\par
\sphinxattableend\end{savenotes}

\sphinxAtStartPar
Enfin, 28\% des personnes interrogées ont déjà produit des données ouvertes. 8\% d’entre elles ont utilisé au moins une fois des entrepôts généralistes comme Nakala (N=5), Zenodo (N=4) ou Data Paris 8 (N=1). Les autres modalité d’ouverture des données (\hyperref[\detokenize{script_analyse/stat_descriptive/connaissance_science_ouverte:autre-entrepot-fig}]{Fig.\@ \ref{\detokenize{script_analyse/stat_descriptive/connaissance_science_ouverte:autre-entrepot-fig}}}) sont Open Science Framework (N=10), des sites web personnels (N=6), des dépôts “git” (N=3) et des entrepots spécialisés (N=3) liés à la linguistiques : Ortolang et Cocoon.

\begin{figure}[htbp]
\centering
\capstart

\noindent\sphinxincludegraphics{{autre_entrepot_rec}.png}
\caption{Les modes d’ouverture des données}\label{\detokenize{script_analyse/stat_descriptive/connaissance_science_ouverte:autre-entrepot-fig}}\end{figure}


\section{Les besoins d’accompagnement à la Science Ouverte}
\label{\detokenize{script_analyse/stat_descriptive/connaissance_science_ouverte:les-besoins-d-accompagnement-a-la-science-ouverte}}
\sphinxAtStartPar
Mis à part l’accompagnement au dépot dans Hal considéré comme inutile par 68\% des répondants, une majorité d’entre eux sont favorables aux accompagnement à la Science Ouverte (\hyperref[\detokenize{script_analyse/stat_descriptive/connaissance_science_ouverte:utilite-accomp-so-fig}]{Fig.\@ \ref{\detokenize{script_analyse/stat_descriptive/connaissance_science_ouverte:utilite-accomp-so-fig}}}). Les licences de diffusion des résultats de recherche est le sujet pour lequel les répondants ont manifesté le plus d’intérêt : 87\% considèrent qu’un accompagnement sur ce thème serait “plutôt utile” (52\%) voire “très utile” (35\%). Ensuite, 81\% des participants à l’enquêtre déclarent qu’il serait utile d’être accompagner dans le processus permettant de produire des données respectant les principes FAIR (Facile à trouver, Accessible, Interopérable, et Réutilisable). De façon général, les accompagnements apparaissent d’autant plus utiles qu’ils portent sur des aspects juridiques (licence, RGPD) ou des sujets liés à l’ouverture des données (métadonnées, dépots, principes FAIR).

\begin{figure}[htbp]
\centering
\capstart

\noindent\sphinxincludegraphics{{util_accomp_so}.png}
\caption{Utilité des accompagnements à la Science Ouverte}\label{\detokenize{script_analyse/stat_descriptive/connaissance_science_ouverte:utilite-accomp-so-fig}}\end{figure}


\section{Stockage, sauvegarde et archivage des données}
\label{\detokenize{script_analyse/stat_descriptive/connaissance_science_ouverte:stockage-sauvegarde-et-archivage-des-donnees}}
\sphinxAtStartPar
Le stockage, la sauvegarde et l’archivage renvoient à des temps différents du cycle de la donnée. Le stockage correspond au moment où les données sont exploitées. Au contraire l’archivage intervient après la phase de traitement et d’analyse, dans un moment où les données ne sont plus manipulées. Enfin, la sauvegarde est une mesure visant à garantir la sécurité et l’integrité des données. Il est ainsi recommandé de copier ses données sur trois supports différents. Il est toutefois possible que les différences entre ces trois concepts ne soient pas connues de tous.

\sphinxAtStartPar
En termes de pratiques, en moyenne, les personnes interrogées sauvegardent leur données sur deux supports et en utilisent trois pour le stockage. La clé usb ou le disque dur externe sont les supports les plus utilisés que ce soit pour le stockage, la sauvegarde ou l’achivage. Concernant les clouds, ce sont les services commerciaux (Google, Icloud) qui arrivent en tête. 49\% des répondants disent stocker ou sauvegarder leurs documents sur ces services et 42\% les emploient comme espace d’achivage. À l’opposé, seulement 9\% des personnes interrogées recourent au cloud de Paris 8 (5\% si on parle de sauvegarde ou d’archivage).

\begin{figure}[htbp]
\centering
\capstart

\noindent\sphinxincludegraphics{{pratiques_sauvegardes}.png}
\caption{Les lieux de stockage, de sauvegarde et d’archivage des données}\label{\detokenize{script_analyse/stat_descriptive/connaissance_science_ouverte:q33-save-data}}\end{figure}

\sphinxstepscope


\chapter{Production et analyse des données}
\label{\detokenize{script_analyse/stat_descriptive/production_et_analyse_donnees:production-et-analyse-des-donnees}}\label{\detokenize{script_analyse/stat_descriptive/production_et_analyse_donnees::doc}}

\section{Types de données produites}
\label{\detokenize{script_analyse/stat_descriptive/production_et_analyse_donnees:types-de-donnees-produites}}
\sphinxAtStartPar
Les personnes interrogées travaillant exclusivement avec des données dites “quantitatives” sont minoritaires (\hyperref[\detokenize{script_analyse/stat_descriptive/production_et_analyse_donnees:quali-or-quanti}]{Table \ref{\detokenize{script_analyse/stat_descriptive/production_et_analyse_donnees:quali-or-quanti}}}), près de la moitié s’appuient essentiellement sur des données qualitatives et environ un tiers manipulent les deux types de données.


\begin{savenotes}\sphinxattablestart
\sphinxthistablewithglobalstyle
\centering
\sphinxcapstartof{table}
\sphinxthecaptionisattop
\sphinxcaption{Données qualitatives ou quantitatives ?}\label{\detokenize{script_analyse/stat_descriptive/production_et_analyse_donnees:quali-or-quanti}}
\sphinxaftertopcaption
\begin{tabulary}{\linewidth}[t]{TTT}
\sphinxtoprule

\sphinxAtStartPar

&\sphinxstyletheadfamily 
\sphinxAtStartPar
nb
&\sphinxstyletheadfamily 
\sphinxAtStartPar
freq
\\
\sphinxmidrule
\sphinxtableatstartofbodyhook
\sphinxAtStartPar
Des données qualitatives
&
\sphinxAtStartPar
57
&
\sphinxAtStartPar
47,9
\\
\sphinxhline
\sphinxAtStartPar
Des données quantitatives
&
\sphinxAtStartPar
19
&
\sphinxAtStartPar
16,0
\\
\sphinxhline
\sphinxAtStartPar
Les deux
&
\sphinxAtStartPar
43
&
\sphinxAtStartPar
36,1
\\
\sphinxhline
\sphinxAtStartPar
Total
&
\sphinxAtStartPar
119
&
\sphinxAtStartPar
100,0
\\
\sphinxbottomrule
\end{tabulary}
\sphinxtableafterendhook\par
\sphinxattableend\end{savenotes}

\sphinxAtStartPar
Les quatre types de données les plus fréquemment produites (\hyperref[\detokenize{script_analyse/stat_descriptive/production_et_analyse_donnees:type-donnee-quali}]{Fig.\@ \ref{\detokenize{script_analyse/stat_descriptive/production_et_analyse_donnees:type-donnee-quali}}})sont les entretiens, les données textuelles, les données d’enquête et les données expérimentales ou d’observation. On notera toutefois le caractère ambigü de certaines catégories. Ainsi la notion d‘“enquête” est largement utilisé en sciences sociales pour désigner la partie empirique des recherche et peut renvoyer aussi bien aux méthodes d’enquêtes ethnographiques qu’aux méthodes d’enquêtes par questionnaire. De même, les données d’observation peuvent se comprendre au sens de l’anthropologue qui observe les interactions au sein d’un groupe humain  ou de l’ornithologue qui compte les espèces d’oiseau. Parmi les “autres types de données”, une personne a d’ailleurs précisé qu’il s’agissait d “observations in situ non expérimentales”. Tandis qu’une autre ne se reconnait pas dans la notion de “donnée”.

\begin{figure}[htbp]
\centering
\capstart

\noindent\sphinxincludegraphics{{type_donnee_quali}.png}
\caption{Les types de données traitées}\label{\detokenize{script_analyse/stat_descriptive/production_et_analyse_donnees:type-donnee-quali}}\end{figure}


\section{Autres types de données}
\label{\detokenize{script_analyse/stat_descriptive/production_et_analyse_donnees:autres-types-de-donnees}}

\section{Outils de traitement}
\label{\detokenize{script_analyse/stat_descriptive/production_et_analyse_donnees:outils-de-traitement}}
\sphinxAtStartPar
Le graphique ci\sphinxhyphen{}dessous (\hyperref[\detokenize{script_analyse/stat_descriptive/production_et_analyse_donnees:logiciel-traitement-use-by-2people}]{Fig.\@ \ref{\detokenize{script_analyse/stat_descriptive/production_et_analyse_donnees:logiciel-traitement-use-by-2people}}}) donne à voir les logiciels qui sont cités par plus de deux répondants. Il regroupe les questions 5 et 5.1 sur les logiciels et langages utilisés pour le traitement des données. La question 5 proposait 6 outils : Excel, R, python, SAS, Stata et Julia. Pour chacun de ces outils, les personnes interrogées devaient répondre si elles les utilisaient ou non. Avec la question 5.1, elles pouvaient ensuite compléter leur réponse en donnant le nom d’autres outils informatiques. Les 36 personnes qui ont répondu à la question 5.1 ont ainsi cité 26 autres logiciels et langages, ce qui fait au total 48  outils informatiques différents. Pour faciliter la lecture, nous les avons regroupés par “familles”. On retrouve ainsi :
\begin{itemize}
\item {} 
\sphinxAtStartPar
les “logiciels de statistiques” comme SPSS, Jamovi, Stata ;

\item {} 
\sphinxAtStartPar
les langages de “scripts et de programmations” comme R, Python, Matlab ou C++. Probablement qu’une subdivision de cette catégorie serait justifiée afin de distinguer les langages généralement employés pour l’analyse de données (R, Julia, Matlab, Python) et les langages de programmation comme C++ ;

\item {} 
\sphinxAtStartPar
les outils de gestion de bibliographie (Zotero, Endnote) ;

\item {} 
\sphinxAtStartPar
les logiciels de “tableurs” et les feuilles de calculs : Excel, Libre Calc ;

\item {} 
\sphinxAtStartPar
les langages liés au développement web (css, html, php, javascript) ;

\item {} 
\sphinxAtStartPar
les outils d’analyse textuelle (Iramuteq, ActiveTigger) ;

\item {} 
\sphinxAtStartPar
les “CAQDAS” (Computer\sphinxhyphen{}Assisted Qualitative Data Analysis Software) comme Nvivo, Qualcoder, Atlas.ti

\end{itemize}

\sphinxAtStartPar
En termes d’utilisateurs, Excel est l’outil le plus répendu. Deux tiers des répondants (N=79) y recourent pour traiter leurs données. Ce sont ensuite les langage R (N=38) et Python (N=23) qui sont les plus utilisés. Jamovi (N=9) est le plus cité des logiciels d’analyse statistique, comme Nvivo (N=6) pour les logiciels CAQDAS. On note par ailleurs qu’une peronne sur cinq dit n’utiliser aucun outil informatique pour traiter les données.

\begin{figure}[htbp]
\centering
\capstart

\noindent\sphinxincludegraphics{{logiciel_traitement_use_by_2people}.png}
\caption{Logiciels cités par plus de deux répondants}\label{\detokenize{script_analyse/stat_descriptive/production_et_analyse_donnees:logiciel-traitement-use-by-2people}}\end{figure}


\section{Développement et dépôt de codes ou de logiciels}
\label{\detokenize{script_analyse/stat_descriptive/production_et_analyse_donnees:developpement-et-depot-de-codes-ou-de-logiciels}}
\sphinxAtStartPar
Parallèlement  à l’utilisation d’outils informatiques pour le traitement des données, près d’un quart des répondants ont déjà été amenés à développer ou faire développer du code ou des logiciels dans le cadre de leurs recherches, dont deux n’ont jamais déposé leur code ou logiciel sur un entrepot dédié.

\sphinxAtStartPar
Si le dépôt de code sur Github et Gitlab restent les solutions les plus usitées, 6 personnes déposent leus codes et logiciels ailleurs :
\begin{itemize}
\item {} 
\sphinxAtStartPar
“Sur les dépôts SSC des commandes Stata”;

\item {} 
\sphinxAtStartPar
“sur des repositories personnels”;

\item {} 
\sphinxAtStartPar
“sur Hal”;

\item {} 
\sphinxAtStartPar
“sur Ortolang et Huma\sphinxhyphen{}Num”,

\item {} 
\sphinxAtStartPar
“Sur des instances auto\sphinxhyphen{}hébergées de Gitlab”

\item {} 
\sphinxAtStartPar
ou “sur l’INPI”.

\end{itemize}

\sphinxAtStartPar
Sur les 32 personnes ayant répondu “Oui” à la question “êtes\sphinxhyphen{}vous amené.e à développer ou faire développer des logiciels/du code spécialisé(s) ?”, 31 ont donnée des indications plus ou moins précises sur ce qu’elles développent ou font développer en répondant à la question 16 (\hyperref[\detokenize{script_analyse/stat_descriptive/production_et_analyse_donnees:dev-code-tab}]{Table \ref{\detokenize{script_analyse/stat_descriptive/production_et_analyse_donnees:dev-code-tab}}}). Nous avons analysé les réponses manuellement. On a ainsi 8 personnes qui disent développer des “logiciels variés”, des “algorithmes” ou encore des “plateformes, programmes logiciels, œuvres” sans données d’autres précisions. Les chercheurs qui manipulent du code afin de traiter et d’analyser les données constituent le groupe le plus importants (N=14). On observe enfin un petit groupe de personnes (N=7) construisant des outils liés à la collecte de données d’expérience via des jeux ou le logiciel Psytoolkit.


\begin{savenotes}\sphinxattablestart
\sphinxthistablewithglobalstyle
\centering
\sphinxcapstartof{table}
\sphinxthecaptionisattop
\sphinxcaption{Types de code développé}\label{\detokenize{script_analyse/stat_descriptive/production_et_analyse_donnees:dev-code-tab}}
\sphinxaftertopcaption
\begin{tabulary}{\linewidth}[t]{TTTT}
\sphinxtoprule
\sphinxstyletheadfamily 
\sphinxAtStartPar
Type de développement
&\sphinxstyletheadfamily 
\sphinxAtStartPar
nb
&\sphinxstyletheadfamily 
\sphinxAtStartPar
total
&\sphinxstyletheadfamily 
\sphinxAtStartPar
freq
\\
\sphinxmidrule
\sphinxtableatstartofbodyhook
\sphinxAtStartPar
traitement et analyse de données
&
\sphinxAtStartPar
15
&
\sphinxAtStartPar
31
&
\sphinxAtStartPar
48,4
\\
\sphinxhline
\sphinxAtStartPar
programmes et logiciels variés
&
\sphinxAtStartPar
8
&
\sphinxAtStartPar
31
&
\sphinxAtStartPar
25,8
\\
\sphinxhline
\sphinxAtStartPar
collecte de données
&
\sphinxAtStartPar
7
&
\sphinxAtStartPar
31
&
\sphinxAtStartPar
22,6
\\
\sphinxhline
\sphinxAtStartPar
diffusion
&
\sphinxAtStartPar
2
&
\sphinxAtStartPar
31
&
\sphinxAtStartPar
6,5
\\
\sphinxbottomrule
\end{tabulary}
\sphinxtableafterendhook\par
\sphinxattableend\end{savenotes}


\section{Accompagnements et formations méthodologiques}
\label{\detokenize{script_analyse/stat_descriptive/production_et_analyse_donnees:accompagnements-et-formations-methodologiques}}

\subsection{Méthodes}
\label{\detokenize{script_analyse/stat_descriptive/production_et_analyse_donnees:methodes}}

\subsection{Outils}
\label{\detokenize{script_analyse/stat_descriptive/production_et_analyse_donnees:outils}}
\sphinxAtStartPar
Les questions concernant les besoins d’accompagnement à différents types de logiciels open source a été posé deux fois :
\begin{itemize}
\item {} 
\sphinxAtStartPar
question 7, partie “Gestion et traitement des données”

\item {} 
\sphinxAtStartPar
la question 3 de la partie “Analyse des données”.

\end{itemize}

\sphinxAtStartPar
La comparaison des tris à plat montre que certaines personnes ont répondu différrement. La différence reste néanmoins marginale. Par exemple, la première fois que la question a été posée, 39 personnes ressentaient le besoin d’être accompagné pour utiliser un tableur contre 30 la seconde fois que la question a été posée (\hyperref[\detokenize{script_analyse/stat_descriptive/production_et_analyse_donnees:tableur-comp-1-2}]{Table \ref{\detokenize{script_analyse/stat_descriptive/production_et_analyse_donnees:tableur-comp-1-2}}}).


\begin{savenotes}\sphinxattablestart
\sphinxthistablewithglobalstyle
\centering
\sphinxcapstartof{table}
\sphinxthecaptionisattop
\sphinxcaption{Besoin d’accompagnement pour l’utilisation de tableur: comparaison des réponses 1 et 2}\label{\detokenize{script_analyse/stat_descriptive/production_et_analyse_donnees:tableur-comp-1-2}}
\sphinxaftertopcaption
\begin{tabulary}{\linewidth}[t]{TTT}
\sphinxtoprule
\sphinxstyletheadfamily 
\sphinxAtStartPar
Tableur
&\sphinxstyletheadfamily 
\sphinxAtStartPar
Réponse 1
&\sphinxstyletheadfamily 
\sphinxAtStartPar
Réponse 2
\\
\sphinxmidrule
\sphinxtableatstartofbodyhook
\sphinxAtStartPar
Non
&
\sphinxAtStartPar
80
&
\sphinxAtStartPar
89
\\
\sphinxhline
\sphinxAtStartPar
Oui, avancé
&
\sphinxAtStartPar
23
&
\sphinxAtStartPar
16
\\
\sphinxhline
\sphinxAtStartPar
Oui, débutant
&
\sphinxAtStartPar
16
&
\sphinxAtStartPar
14
\\
\sphinxhline
\sphinxAtStartPar
Total
&
\sphinxAtStartPar
119
&
\sphinxAtStartPar
119
\\
\sphinxbottomrule
\end{tabulary}
\sphinxtableafterendhook\par
\sphinxattableend\end{savenotes}

\sphinxAtStartPar
Dans l’ensemble, entre un quart et un tiers des répondants estime avoir besoin d’un accompagnement pour les différents logiciels cités (\hyperref[\detokenize{script_analyse/stat_descriptive/production_et_analyse_donnees:help-outil}]{Fig.\@ \ref{\detokenize{script_analyse/stat_descriptive/production_et_analyse_donnees:help-outil}}}). Les quatre premier types de logiciel demandés sont :
\begin{enumerate}
\sphinxsetlistlabels{\arabic}{enumi}{enumii}{}{.}%
\item {} 
\sphinxAtStartPar
les logiciels d’analyse textuelle. 48\% des répondants disent ressentir le besoin d’un accompagnement (35\% à un niveau débutant et 13\% à un niveau avancé). À noter que les logiciels cités renvoient à des approches et des méthodes d’analyse différentes.

\item {} 
\sphinxAtStartPar
les logiciel d’analyse statistique (40\%)

\item {} 
\sphinxAtStartPar
les logiciels d’analyse de données qualitatives comme Qualcoder (environ 37\%)

\item {} 
\sphinxAtStartPar
le langage python (environ 35\%)

\end{enumerate}

\sphinxAtStartPar
Ensuite, on observe des demandes d’accompagnement pour les logiciels concernant des méthodes, des données ou des langages plus spécifiques comme les outils d’analyse de réseau de cartographie ou de gestion de bases de données. En bas de classement, nous trouvons enfin les logiciels d’océrisation et de système d’information géographique.

\sphinxAtStartPar
Dans le cas des besoins en accompagnement à la prise en main de logiciels propriétaires, 5 “familles” de logiciels étaient proposaient :
\begin{itemize}
\item {} 
\sphinxAtStartPar
Les tableurs représentés par Excel

\item {} 
\sphinxAtStartPar
les logiciels de traitements statistiques comme SPSS, Stata, SAS

\item {} 
\sphinxAtStartPar
les logiciels de traitement des données textuelles avec Atlas.ti

\item {} 
\sphinxAtStartPar
les systèmes d’information géographique via ArcGis

\item {} 
\sphinxAtStartPar
les logiciels d’analyse des données qualitative (CAQDAS, Nvivo)

\end{itemize}

\sphinxAtStartPar
Le taux de personnes exprimant un besoin d’accompagnement est globablement plus faible que pour les outils open source. Excel est l’outil qui obtient le taux le plus élevé : environ 27\% des personnes disent ressentir le besoin d’être accompagnées. On retrouve une proportion relativement similaire à celles des outils de tableur en open source. On constate par contre un écart de 20 personnes entre celles qui sont intéressées par un accompagnement aux logiciels CAQDAS open source et Nvivo (\hyperref[\detokenize{script_analyse/stat_descriptive/production_et_analyse_donnees:q20-19-nvivo}]{Table \ref{\detokenize{script_analyse/stat_descriptive/production_et_analyse_donnees:q20-19-nvivo}}}).

\begin{figure}[htbp]
\centering
\capstart

\noindent\sphinxincludegraphics{{q16_q19_help_outil}.png}
\caption{Besoin d’accompagnement aux outils open source}\label{\detokenize{script_analyse/stat_descriptive/production_et_analyse_donnees:help-outil}}\end{figure}

\begin{figure}[htbp]
\centering
\capstart

\noindent\sphinxincludegraphics{{script_analyse/stat_descriptive/viz/9_help_outil_proprio}.png)}
\caption{Besoin d’accompagnement aux outils propriétaires}\label{\detokenize{script_analyse/stat_descriptive/production_et_analyse_donnees:help-outil-proprio}}\end{figure}

\sphinxAtStartPar
Il semblerait donc que les chercheurs de Paris 8 sont d’avantage intéressés par des outils open source plutôt que propriétaires. Il faut toutefois rester prudent sur cette interprétation des écarts observés dans la mesure où ils peuvent être également lié à une fatigue des répondants, ces derniers choisissant “non” par facilité. Il est possible aussi que la question sur les logiciels open source apparaissent plus générale : en répondant “oui je ressens un besoin d’accompagnement à un logiciel d’analyse textuelle”, les personnes visent autant un outil en particulier que la méthode en elle\sphinxhyphen{}même (indépendamment du logiciel utilisé).

\sphinxAtStartPar
Une manière de valider l’idée d’une plus grande affinité des chercheurs de Paris 8 pour les logiciels libres serait de comparer ces résultats avec ceux du sondage de la DSIN sur le type de suite bureautique (microsoft, google, outils libres) péblicitée par le personnel de l’université.


\begin{savenotes}\sphinxattablestart
\sphinxthistablewithglobalstyle
\centering
\sphinxcapstartof{table}
\sphinxthecaptionisattop
\sphinxcaption{Comparaison des besoins en accompagnement à Nvivo et au logiciels CAQDAS open source}\label{\detokenize{script_analyse/stat_descriptive/production_et_analyse_donnees:q20-19-nvivo}}
\sphinxaftertopcaption
\begin{tabulary}{\linewidth}[t]{TTTT}
\sphinxtoprule

\sphinxAtStartPar

&\sphinxstyletheadfamily 
\sphinxAtStartPar
Non, je ne ressens pas le besoin d’un accompagnement à un outil CAQDAS open source
&\sphinxstyletheadfamily 
\sphinxAtStartPar
je ressens le besoin d’un accompagnement à un outil CAQDAS open source
&\sphinxstyletheadfamily 
\sphinxAtStartPar
Total
\\
\sphinxmidrule
\sphinxtableatstartofbodyhook
\sphinxAtStartPar
Non, je ne ressens pas le besoin d’un accompagnement à Nvivo
&
\sphinxAtStartPar
66
&
\sphinxAtStartPar
26
&
\sphinxAtStartPar
92
\\
\sphinxhline
\sphinxAtStartPar
Oui, je ressens le besoin d’un accompagnement à Nvivo
&
\sphinxAtStartPar
6
&
\sphinxAtStartPar
21
&
\sphinxAtStartPar
27
\\
\sphinxhline
\sphinxAtStartPar
Total
&
\sphinxAtStartPar
72
&
\sphinxAtStartPar
47
&
\sphinxAtStartPar
119
\\
\sphinxbottomrule
\end{tabulary}
\sphinxtableafterendhook\par
\sphinxattableend\end{savenotes}

\sphinxAtStartPar
Enfin, une dizaine répondants ont fait des suggestion d’outils open source ou propriétaires concernant lesquels ils aimeraient être formés. Ces outils s’inscrivent dans les différentes catégories proposées comme “tropy” (analyses de données qualitatives), jamovi (analyse statistique), les bases de données “orientées graphe” (système de gestion de base de données) ou des outils de traitement du son, des textes ou des images.


\begin{savenotes}\sphinxattablestart
\sphinxthistablewithglobalstyle
\centering
\sphinxcapstartof{table}
\sphinxthecaptionisattop
\sphinxcaption{Suggestion d’accompagnement à des outils open source ou propriétaire}\label{\detokenize{script_analyse/stat_descriptive/production_et_analyse_donnees:suggestion-software}}
\sphinxaftertopcaption
\begin{tabulary}{\linewidth}[t]{TT}
\sphinxtoprule
\sphinxstyletheadfamily 
\sphinxAtStartPar
Autres logiciels open source
&\sphinxstyletheadfamily 
\sphinxAtStartPar
Autres logiciels propriétaires
\\
\sphinxmidrule
\sphinxtableatstartofbodyhook
\sphinxAtStartPar
logiciel d’analyse des gestes, de l’intonation etc.
&
\sphinxAtStartPar
Qualtrics, Jamovi
\\
\sphinxhline
\sphinxAtStartPar
wikidata
&
\sphinxAtStartPar
Photoshop
\\
\sphinxhline
\sphinxAtStartPar
Logiciel de reconnaissance de texte manuscrits type Transkribus
&
\sphinxAtStartPar
Jamovi
\\
\sphinxhline
\sphinxAtStartPar
modèle linéaire mixte (LMM)
&
\sphinxAtStartPar

\\
\sphinxhline
\sphinxAtStartPar
Bases de données en graphe éventuellement
&
\sphinxAtStartPar

\\
\sphinxhline
\sphinxAtStartPar
Logiciel de traitement d’entretiens et de vidéos
&
\sphinxAtStartPar

\\
\sphinxhline
\sphinxAtStartPar
Anaconda (python), Tropes…
&
\sphinxAtStartPar

\\
\sphinxhline
\sphinxAtStartPar
Tropy
&
\sphinxAtStartPar

\\
\sphinxbottomrule
\end{tabulary}
\sphinxtableafterendhook\par
\sphinxattableend\end{savenotes}


\subsection{Outils propriétaires}
\label{\detokenize{script_analyse/stat_descriptive/production_et_analyse_donnees:outils-proprietaires}}

\subsection{Services à développer}
\label{\detokenize{script_analyse/stat_descriptive/production_et_analyse_donnees:services-a-developper}}
\sphinxAtStartPar
Enfin, les services à développer proposés dans le questionnaire ont rencontré un plébiscite dans la mesure où la majorité des personnes les considèrent comme utiles. Le service ayant le moins de succès concerne l’accès à une SD\sphinxhyphen{}Box et au CASD (Centre d’accès sécurisé aux données). Près de 63\% le jugent utile contre 84\% pour les ressources documentaires ou 80\% pour les permanences d’accueil. Parmi les 44 personnes non\sphinxhyphen{}intéressées par ce dernier service (36\%), il est possible qu’une partie d’entre elles ne connaissent pas les services du CASD et la “SD\sphinxhyphen{}Box”.
\begin{quote}

\sphinxAtStartPar
La SD\sphinxhyphen{}BOX “La SD\sphinxhyphen{}Box est le seul moyen d’accès à l’infrastructure centrale du CASD et à l’environnement de travail de l’utilisateur. C’est un dispositif technique conçu et mis au point par le CASD pour répondre aux besoins et contraintes des utilisateurs et déposants de données”.
\end{quote}

\begin{figure}[htbp]
\centering
\capstart

\noindent\sphinxincludegraphics{{11_ressources}.png}
\caption{Les services à développer}\label{\detokenize{script_analyse/stat_descriptive/production_et_analyse_donnees:service-to-develop}}\end{figure}

\sphinxAtStartPar
Par ailleurs, une quinzaine de personnes ont répondu à la question ouverte : “d’autres services vous paraitraient\sphinxhyphen{}ils utiles ?”. Plusieurs réponse portent sur les ressources documentaires de la bibliothèque qu’ils jugent “pauvres” au regard d’autre centres ou davantage orientées vers la formation que la recherche.
\begin{quote}

\sphinxAtStartPar
“Quand je parle de documentation je parle d’ouvrages pour la recherche la bibliothèque étant très loin d’être au niveau sur cet aspect même si elle a un bon.fonds pour la formation.”
\end{quote}
\begin{quote}

\sphinxAtStartPar
“Il est certain que les ressources documentaires de la bibliothèque sont très pauvres en ce qui concerne la recherche et ne peuvent être utilisées, à mon sens, par des chercheurs d’aujourd’hui.”
\end{quote}

\sphinxAtStartPar
Les autres commentaires portent sur l’accès à des bureaux nominatifs et équipés informatiquement, “la possibilité de créer des accès wifi pour les stagiaires et les nouveaux collègues dans un délai raisonnable”, le développement de services liés à la sécurité des données (accès à un serveur chiffré et des lieux de stockages sécurisés) ou de formation sur l’éthique de la recherche.
\begin{quote}

\sphinxAtStartPar
“une formation des chercheurs à des obligations/engagements éthiques ; le comité d’éthique de P8 est récent mais son existence est peu connu ;”
\end{quote}

\sphinxstepscope


\chapter{Le TGIR Huma\sphinxhyphen{}num et ses outils}
\label{\detokenize{script_analyse/stat_descriptive/Outils_humanum:le-tgir-huma-num-et-ses-outils}}\label{\detokenize{script_analyse/stat_descriptive/Outils_humanum::doc}}
\sphinxAtStartPar
La majorité des répondants (55\%, \hyperref[\detokenize{script_analyse/stat_descriptive/Outils_humanum:humanum}]{Table \ref{\detokenize{script_analyse/stat_descriptive/Outils_humanum:humanum}}}) disent ne pas connaître l’infrastructure Huma\sphinxhyphen{}Num. Concernant les services, {\hyperref[\detokenize{script_analyse/stat_descriptive/glossaire:term-ShareDocs}]{\sphinxtermref{\DUrole{xref,std,std-term}{Sharedocs}}}} est celui qui est le plus ancré dans les usages avec 38\% des personnes interrogées qui disent l’avoir déjà utilisé. Viennent ensuit {\hyperref[\detokenize{script_analyse/stat_descriptive/glossaire:term-Isidore}]{\sphinxtermref{\DUrole{xref,std,std-term}{ISIDORE}}}} (18\%),  {\hyperref[\detokenize{script_analyse/stat_descriptive/glossaire:term-Gitlab}]{\sphinxtermref{\DUrole{xref,std,std-term}{Gitlab}}}} (13\%) et {\hyperref[\detokenize{script_analyse/stat_descriptive/glossaire:term-Nakala}]{\sphinxtermref{\DUrole{xref,std,std-term}{Nakala}}}} (13\%). Les autres services sont encore plus confidentiels. {\hyperref[\detokenize{script_analyse/stat_descriptive/glossaire:term-Stylo}]{\sphinxtermref{\DUrole{xref,std,std-term}{Stylo}}}}, un outil de rédaction adapté à l’écriture en sciences sociales, n’a été utilisé que par 4\% des répondants.


\begin{savenotes}\sphinxattablestart
\sphinxthistablewithglobalstyle
\centering
\sphinxcapstartof{table}
\sphinxthecaptionisattop
\sphinxcaption{Connaissance de la TGIR Huma\sphinxhyphen{}Num}\label{\detokenize{script_analyse/stat_descriptive/Outils_humanum:humanum}}
\sphinxaftertopcaption
\begin{tabulary}{\linewidth}[t]{TTT}
\sphinxtoprule

\sphinxAtStartPar

&\sphinxstyletheadfamily 
\sphinxAtStartPar
nb
&\sphinxstyletheadfamily 
\sphinxAtStartPar
freq
\\
\sphinxmidrule
\sphinxtableatstartofbodyhook
\sphinxAtStartPar
Non
&
\sphinxAtStartPar
66
&
\sphinxAtStartPar
55.5
\\
\sphinxhline
\sphinxAtStartPar
Oui
&
\sphinxAtStartPar
53
&
\sphinxAtStartPar
44.5
\\
\sphinxhline
\sphinxAtStartPar
Total
&
\sphinxAtStartPar
119
&
\sphinxAtStartPar
100.0
\\
\sphinxbottomrule
\end{tabulary}
\sphinxtableafterendhook\par
\sphinxattableend\end{savenotes}

\sphinxAtStartPar
Sur l’ensemble des services Huma\sphinxhyphen{}Num, une majorité de répondants ne ressent pas le besoin d’un accompagnement. Les accompagnements les plus demandés portent sur l’hébergement web (39\%), Nakala (39\%), Sharedocs (36\%) et Isidore (36\%).

\begin{figure}[htbp]
\centering
\capstart

\noindent\sphinxincludegraphics{{q13_q14_outil_humanum}.png}
\caption{L’utilisation des services Huma\sphinxhyphen{}Num}\label{\detokenize{script_analyse/stat_descriptive/Outils_humanum:q13-q14-humanum}}\end{figure}

\begin{figure}[htbp]
\centering
\capstart

\noindent\sphinxincludegraphics{{q15_help_outil_humanum}.png}
\caption{Les besoins en accompagnement aux outils Huma\sphinxhyphen{}Num}\label{\detokenize{script_analyse/stat_descriptive/Outils_humanum:q15-help-outil-hn}}\end{figure}

\sphinxAtStartPar
Le croisement de la question “Connaissez\sphinxhyphen{}vous la TGIR Huma\sphinxhyphen{}num” et celles sur les outils utilisés montre que ceux\sphinxhyphen{}ci ne sont pas toujours associés à Huma\sphinxhyphen{}Num. Si toutes les personnes utilisant Nakala connaissent Huma\sphinxhyphen{}Num (\hyperref[\detokenize{script_analyse/stat_descriptive/Outils_humanum:nakala-x-humanum}]{Table \ref{\detokenize{script_analyse/stat_descriptive/Outils_humanum:nakala-x-humanum}}}), ce n’est pas le cas pour ShareDocs (\hyperref[\detokenize{script_analyse/stat_descriptive/Outils_humanum:sharedocs-x-humanum}]{Table \ref{\detokenize{script_analyse/stat_descriptive/Outils_humanum:sharedocs-x-humanum}}}) : une dizaine d’utilisateur de ce service ont répondu qu’elles ne connaissez pas la TGIR. On dénombre également 7 personnes connaissant Huma\sphinxhyphen{}Num mais n’ayant utilisé aucun des services de l’infrastructure.


\begin{savenotes}\sphinxattablestart
\sphinxthistablewithglobalstyle
\centering
\sphinxcapstartof{table}
\sphinxthecaptionisattop
\sphinxcaption{Utilisation de Nakala et connaissance d’Huma\sphinxhyphen{}Num}\label{\detokenize{script_analyse/stat_descriptive/Outils_humanum:nakala-x-humanum}}
\sphinxaftertopcaption
\begin{tabulary}{\linewidth}[t]{TTTT}
\sphinxtoprule

\sphinxAtStartPar

&\sphinxstyletheadfamily 
\sphinxAtStartPar
Non, je n’ai jamais utilisé Nakala
&\sphinxstyletheadfamily 
\sphinxAtStartPar
Oui, j’ai déjà utilisé Nakala
&\sphinxstyletheadfamily 
\sphinxAtStartPar
Total
\\
\sphinxmidrule
\sphinxtableatstartofbodyhook
\sphinxAtStartPar
Non, je ne connais pas Huma\sphinxhyphen{}Num
&
\sphinxAtStartPar
66
&
\sphinxAtStartPar
0
&
\sphinxAtStartPar
66
\\
\sphinxhline
\sphinxAtStartPar
Oui, je connais Huma\sphinxhyphen{}Num
&
\sphinxAtStartPar
38
&
\sphinxAtStartPar
15
&
\sphinxAtStartPar
53
\\
\sphinxhline
\sphinxAtStartPar
Total
&
\sphinxAtStartPar
104
&
\sphinxAtStartPar
15
&
\sphinxAtStartPar
119
\\
\sphinxbottomrule
\end{tabulary}
\sphinxtableafterendhook\par
\sphinxattableend\end{savenotes}


\begin{savenotes}\sphinxattablestart
\sphinxthistablewithglobalstyle
\centering
\sphinxcapstartof{table}
\sphinxthecaptionisattop
\sphinxcaption{Utilisation de ShareDocs et connaissance d’Huma\sphinxhyphen{}Num}\label{\detokenize{script_analyse/stat_descriptive/Outils_humanum:sharedocs-x-humanum}}
\sphinxaftertopcaption
\begin{tabulary}{\linewidth}[t]{TTTT}
\sphinxtoprule

\sphinxAtStartPar

&\sphinxstyletheadfamily 
\sphinxAtStartPar
Non, je n’ai jamais utilisé ShareDocs
&\sphinxstyletheadfamily 
\sphinxAtStartPar
Oui, j’ai déjà utilisé ShareDocs
&\sphinxstyletheadfamily 
\sphinxAtStartPar
All
\\
\sphinxmidrule
\sphinxtableatstartofbodyhook
\sphinxAtStartPar
Non, je ne connais pas Huma\sphinxhyphen{}Num
&
\sphinxAtStartPar
56
&
\sphinxAtStartPar
10
&
\sphinxAtStartPar
66
\\
\sphinxhline
\sphinxAtStartPar
Oui, je connais Huma\sphinxhyphen{}Num
&
\sphinxAtStartPar
18
&
\sphinxAtStartPar
35
&
\sphinxAtStartPar
53
\\
\sphinxhline
\sphinxAtStartPar
Total
&
\sphinxAtStartPar
74
&
\sphinxAtStartPar
45
&
\sphinxAtStartPar
119
\\
\sphinxbottomrule
\end{tabulary}
\sphinxtableafterendhook\par
\sphinxattableend\end{savenotes}

\sphinxstepscope


\chapter{Profils des enquêtés}
\label{\detokenize{script_analyse/stat_descriptive/profil_repondant:profils-des-enquetes}}\label{\detokenize{script_analyse/stat_descriptive/profil_repondant::doc}}

\section{Les pratiques de travail}
\label{\detokenize{script_analyse/stat_descriptive/profil_repondant:les-pratiques-de-travail}}
\sphinxAtStartPar
La première question interrogeant les pratiques de travail des enquêtés porte sur le lieu habituel de travail en classant par ordre de priorité les options suivantes :
\begin{itemize}
\item {} 
\sphinxAtStartPar
“À mon domicile”

\item {} 
\sphinxAtStartPar
“Dans mon bureau sur site”

\item {} 
\sphinxAtStartPar
“Autre”

\end{itemize}

\sphinxAtStartPar
90\% des personnes interrogés disent travailler le plus souvent à leur domicile (\hyperref[\detokenize{script_analyse/stat_descriptive/profil_repondant:env-travail}]{Table \ref{\detokenize{script_analyse/stat_descriptive/profil_repondant:env-travail}}}). C’est aussi le lieu le plus souvent cité en premier. 71\% des répondants travaillent ainsi prioritairement à leur domicile. Tandis que le bureau est en moyenne classé deuxième dans l’ordre des priorités.

\sphinxAtStartPar
Parmis les autres lieux cités, on trouve tout d’abord les bibliothèques et les autres sites (CNRS, Campus Condorcet, partenaires de recherche). On note également que plusieurs répondants ont indiqué qu’ils travaillaient un peu partout, là où ils pouvaient faute de bureaux disponibles sur le campus.
\begin{quote}

\sphinxAtStartPar
“Où je peux, puisque nous n’avon pas de bureaux”
\end{quote}
\begin{quote}

\sphinxAtStartPar
“bureau de collègues dans d’autres universités”
\end{quote}
\begin{quote}

\sphinxAtStartPar
“Bureau MSH OU BIBLIOTHÈQUE MAIS JE MANQUE CRUELLEMENT D’UN ESPACE DE TRAVAIL”
\end{quote}
\begin{quote}

\sphinxAtStartPar
“où je peux, je n’ai pas de bureau à la fac (15 m2 pour 6 titulaires)”
\end{quote}

\sphinxAtStartPar
Ces réactions recoupent les résultats observés concernant les services à développer où la construction d’espace de travail avaient été choisi par plus de 80\% des répondants (\hyperref[\detokenize{script_analyse/stat_descriptive/production_et_analyse_donnees:service-to-develop}]{Fig.\@ \ref{\detokenize{script_analyse/stat_descriptive/production_et_analyse_donnees:service-to-develop}}}).


\begin{savenotes}\sphinxattablestart
\sphinxthistablewithglobalstyle
\centering
\sphinxcapstartof{table}
\sphinxthecaptionisattop
\sphinxcaption{Où travaillez\sphinxhyphen{}vous le plus souvent ?}\label{\detokenize{script_analyse/stat_descriptive/profil_repondant:env-travail}}
\sphinxaftertopcaption
\begin{tabulary}{\linewidth}[t]{TTTT}
\sphinxtoprule
\sphinxstyletheadfamily 
\sphinxAtStartPar
rang
&\sphinxstyletheadfamily 
\sphinxAtStartPar
A mon domicile
&\sphinxstyletheadfamily 
\sphinxAtStartPar
Dans mon bureau sur site
&\sphinxstyletheadfamily 
\sphinxAtStartPar
Autre, précisez
\\
\sphinxmidrule
\sphinxtableatstartofbodyhook
\sphinxAtStartPar
1
&
\sphinxAtStartPar
85 (71,4\%)
&
\sphinxAtStartPar
29 (24,4\%)
&
\sphinxAtStartPar
5  (4,2\%)
\\
\sphinxhline
\sphinxAtStartPar
2
&
\sphinxAtStartPar
20 (16,8\%)
&
\sphinxAtStartPar
36 (20,3\%)
&
\sphinxAtStartPar
12 (10,1\%)
\\
\sphinxhline
\sphinxAtStartPar
3
&
\sphinxAtStartPar
3  (2,5\%)
&
\sphinxAtStartPar
5  (4,2\%)
&
\sphinxAtStartPar
14 (11,8\%)
\\
\sphinxhline
\sphinxAtStartPar
Total
&
\sphinxAtStartPar
108 (90,8\%)
&
\sphinxAtStartPar
70 (58,8\%)
&
\sphinxAtStartPar
31 (26,1\%)
\\
\sphinxbottomrule
\end{tabulary}
\sphinxtableafterendhook\par
\sphinxattableend\end{savenotes}

\sphinxAtStartPar
Les systèmes d’exploitation les plus répandus sont Apple et Windows. Seules 8 personnes utilisent Linux. En ce qui concerne les navigateurs, Firefox est le plus populaire. Il est cité par 52\% des répondants. Chrome vient ensuite avec 23\% des répondants qui l’utilisent.


\begin{savenotes}\sphinxattablestart
\sphinxthistablewithglobalstyle
\centering
\sphinxcapstartof{table}
\sphinxthecaptionisattop
\sphinxcaption{Les systèmes d’exploitation utilisés}\label{\detokenize{script_analyse/stat_descriptive/profil_repondant:os-systeme}}
\sphinxaftertopcaption
\begin{tabulary}{\linewidth}[t]{TTT}
\sphinxtoprule
\sphinxstyletheadfamily 
\sphinxAtStartPar
q29\_os\_sys
&\sphinxstyletheadfamily 
\sphinxAtStartPar
nb
&\sphinxstyletheadfamily 
\sphinxAtStartPar
freq
\\
\sphinxmidrule
\sphinxtableatstartofbodyhook
\sphinxAtStartPar
Apple
&
\sphinxAtStartPar
53
&
\sphinxAtStartPar
44,5
\\
\sphinxhline
\sphinxAtStartPar
ChromeOS
&
\sphinxAtStartPar
4
&
\sphinxAtStartPar
3,4
\\
\sphinxhline
\sphinxAtStartPar
Linux
&
\sphinxAtStartPar
8
&
\sphinxAtStartPar
6,7
\\
\sphinxhline
\sphinxAtStartPar
Windows
&
\sphinxAtStartPar
54
&
\sphinxAtStartPar
45,4
\\
\sphinxhline
\sphinxAtStartPar
Total
&
\sphinxAtStartPar
119
&
\sphinxAtStartPar
100,0
\\
\sphinxbottomrule
\end{tabulary}
\sphinxtableafterendhook\par
\sphinxattableend\end{savenotes}


\section{Utilisation de l’IA}
\label{\detokenize{script_analyse/stat_descriptive/profil_repondant:utilisation-de-l-ia}}
\sphinxAtStartPar
Enfin, 68 répondants (soit 57\%) disent avoir recourt à l’intelligence artificielle pour leurs recherches, dont 65\% emploient des versions gratuites. La même proportion cite Chat GPT. Il est l’outil le plus utilisé. L’analyse de la question “pourquoi les utilisez\sphinxhyphen{}vous ?” à l’aide de l’algorithme de Topic modeling “BERTopic” {[}\hyperlink{cite.biblio:id3}{Grootendorst, n.d.}{]} fait ressortir six motifs avancés par les personnes enquêtées pour expliquer leur utilisation de l’IA. Celui qui revient le plus souvent est le “gain de temps”, que ce soit pour rédiger des e\sphinxhyphen{}mails, effectuer des tâches administratives ou rechercher de l’information :
\begin{quote}

\sphinxAtStartPar
“Pour gagner du temps sur des procédures routinières ou pour formaliser des courriels ou autres plus rapidement”
\end{quote}
\begin{quote}

\sphinxAtStartPar
“Pour gagner du temps dans la recherche informations”
\end{quote}

\sphinxAtStartPar
Les deux autres principaux motifs sont l’aide à la traduction et le traitement des données à travers la retranscription d’entretiens avec Whisper ou l’édition de code informatique.
\begin{quote}

\sphinxAtStartPar
“relecture de l’anglais principalement”
\end{quote}
\begin{quote}

\sphinxAtStartPar
M’aider à écrire du code de traitement de donnée principalement
\end{quote}
\begin{quote}

\sphinxAtStartPar
analyse et visualisation des données
\end{quote}
\begin{quote}

\sphinxAtStartPar
analyse d’entretiens, recherche dans de multiples entretiens
\end{quote}

\sphinxAtStartPar
L’IA est enfin utilisée pour la relecture et la correction de texte, dans un contexte pédagogique, que ce soit pour préparer des cours ou contrôle son utilisation par les étudiants, ou pour comprendre ce que ces outils font.
\begin{quote}

\sphinxAtStartPar
Relectures, améliorations du manuscrit
\end{quote}
\begin{quote}

\sphinxAtStartPar
Enseignement: Structuration de cours, élaboration de syllabus, idées d’activités pédagogiques
\end{quote}
\begin{quote}

\sphinxAtStartPar
J’essaye de comprendre comment ils fonctionnent
\end{quote}

\sphinxAtStartPar
Bien sûr, ces motifs peuvent apparaître au sein d’une même réponse. Par exemple, un des répondants a écrit : “NoScribe en local pour des \sphinxstylestrong{transcriptions d’entretiens} + deepl pour des \sphinxstylestrong{traductions courtes et ponctuelles}”.


\begin{savenotes}\sphinxattablestart
\sphinxthistablewithglobalstyle
\centering
\sphinxcapstartof{table}
\sphinxthecaptionisattop
\sphinxcaption{Les 5 premiers outils IA les plus cités. Les pourcentages sont établis sur les 65 personnes ayant recourt à l’IA}\label{\detokenize{script_analyse/stat_descriptive/profil_repondant:top-5-ia}}
\sphinxaftertopcaption
\begin{tabulary}{\linewidth}[t]{TTTT}
\sphinxtoprule
\sphinxstyletheadfamily 
\sphinxAtStartPar
Outils IA
&\sphinxstyletheadfamily 
\sphinxAtStartPar
nb
&\sphinxstyletheadfamily 
\sphinxAtStartPar
total
&\sphinxstyletheadfamily 
\sphinxAtStartPar
freq
\\
\sphinxmidrule
\sphinxtableatstartofbodyhook
\sphinxAtStartPar
chatgpt
&
\sphinxAtStartPar
42
&
\sphinxAtStartPar
65
&
\sphinxAtStartPar
64,6
\\
\sphinxhline
\sphinxAtStartPar
claude
&
\sphinxAtStartPar
19
&
\sphinxAtStartPar
65
&
\sphinxAtStartPar
29,2
\\
\sphinxhline
\sphinxAtStartPar
deepl
&
\sphinxAtStartPar
7
&
\sphinxAtStartPar
65
&
\sphinxAtStartPar
10,8
\\
\sphinxhline
\sphinxAtStartPar
gemini
&
\sphinxAtStartPar
7
&
\sphinxAtStartPar
65
&
\sphinxAtStartPar
10,8
\\
\sphinxhline
\sphinxAtStartPar
mistral
&
\sphinxAtStartPar
6
&
\sphinxAtStartPar
65
&
\sphinxAtStartPar
9,2
\\
\sphinxbottomrule
\end{tabulary}
\sphinxtableafterendhook\par
\sphinxattableend\end{savenotes}

\begin{figure}[htbp]
\centering
\capstart

\noindent\sphinxincludegraphics{{raison_ia_bertopic}.png}
\caption{6 raisons d’utiliser l’IA}\label{\detokenize{script_analyse/stat_descriptive/profil_repondant:q37-raison-ia}}\end{figure}

\sphinxstepscope


\chapter{Glossaire}
\label{\detokenize{script_analyse/stat_descriptive/glossaire:glossaire}}\label{\detokenize{script_analyse/stat_descriptive/glossaire::doc}}\begin{description}
\sphinxlineitem{Gitlab\index{Gitlab@\spxentry{Gitlab}|spxpagem}\phantomsection\label{\detokenize{script_analyse/stat_descriptive/glossaire:term-Gitlab}}}
\sphinxAtStartPar
L’IR* Huma\sphinxhyphen{}Num met à disposition des projets de recherche en sciences humaines et sociales (SHS) une instance du logiciel Gitlab. L’instance GitLab d’Huma\sphinxhyphen{}Num permet l’hébergement sécurisé et le partage maîtrisé de fichiers, en particulier des fichiers de code informatique (tous langages) selon le protocole de versionnage git.

\sphinxlineitem{HNurl\index{HNurl@\spxentry{HNurl}|spxpagem}\phantomsection\label{\detokenize{script_analyse/stat_descriptive/glossaire:term-HNurl}}}
\sphinxAtStartPar
\sphinxhref{http://HNurl.fr}{HNurl.fr}, pour Huma\sphinxhyphen{}Num URL, est un service en ligne permettant de créer des URL courtes non pérennes. Il est développé et maintenu par l’IR* Huma\sphinxhyphen{}Num pour l’ensemble de la communauté d’enseignement et de recherche des sciences humaines et sociales. Si la création d’une URL courte dans \sphinxhref{http://HNurl.fr}{HNurl.fr} se fait après une connexion via HumanID, l’utilisation de l’URL courte générée est publique, libre et ouverte (par exemple dans un email ou sur les réseaux sociaux).

\sphinxlineitem{Isidore\index{Isidore@\spxentry{Isidore}|spxpagem}\phantomsection\label{\detokenize{script_analyse/stat_descriptive/glossaire:term-Isidore}}}
\sphinxAtStartPar
Service proposé par Huma\sphinxhyphen{}Num, Isidore est un moteur et un assistant de recherche permettant de découvrir et de trouver des publications, des données numériques et profils de chercheurs et chercheures en sciences humaines et sociales (SHS) venant du monde entier. Il s’inscrit dans le périmètre des moteurs de recherche académiques favorisant la découvrabilité des publications et des données de la recherche. Il est conçu, développé et hébergé en France.

\sphinxlineitem{Kanboard\index{Kanboard@\spxentry{Kanboard}|spxpagem}\phantomsection\label{\detokenize{script_analyse/stat_descriptive/glossaire:term-Kanboard}}}
\sphinxAtStartPar
Kanboard est un logiciel permettant la gestion de projets de recherche fondé sur la méthode organisationnelle “Kanban”. Huma\sphinxhyphen{}Num propose aux projets de recherche en SHS une instance Kanboard hébergé sur ses serveurs en France. Ce service est à disposition dans la limite des ressources dont dispose Huma\sphinxhyphen{}Num.

\sphinxlineitem{Mattermost\index{Mattermost@\spxentry{Mattermost}|spxpagem}\phantomsection\label{\detokenize{script_analyse/stat_descriptive/glossaire:term-Mattermost}}}
\sphinxAtStartPar
Huma\sphinxhyphen{}Num met à disposition une instance du logiciel Matomo qui permet d’analyser le trafic et l’usage d’un site web hébergé sur les infrastructures d’Huma\sphinxhyphen{}Num. Cette alternative aux outils de Google, permet de suivre l’usage d’un site web tout en conservant les données utilisateurs au sein de l’infrastructure Huma\sphinxhyphen{}Num. Matomo fonctionne à l’aide d’un “traceur” (un code HTML à ajouter dans le code de vos sites Web) et il est constitué d’une interface web qui permet de visualiser les usages selon divers indicateurs et d’en configurer des nouveaux.

\sphinxlineitem{Nakala\index{Nakala@\spxentry{Nakala}|spxpagem}\phantomsection\label{\detokenize{script_analyse/stat_descriptive/glossaire:term-Nakala}}}
\sphinxAtStartPar
Nakala est un entrepôt de données de recherche pour les Sciences Humaines et Sociales certifié entrepôt de données de confiance par le Core Trust Seal en 2026 et par le Comité pour la science ouverte depuis 2024.

\sphinxlineitem{ShareDocs\index{ShareDocs@\spxentry{ShareDocs}|spxpagem}\phantomsection\label{\detokenize{script_analyse/stat_descriptive/glossaire:term-ShareDocs}}}
\sphinxAtStartPar
ShareDocs est un gestionnaire de fichiers mis en oeuvre par l’IR* Huma\sphinxhyphen{}Num, sur ses propres serveurs, pouvant être utilisé via un navigateur web, un client WebDAV ou un logiciel de synchronisation de fichiers.

\sphinxlineitem{Stylo\index{Stylo@\spxentry{Stylo}|spxpagem}\phantomsection\label{\detokenize{script_analyse/stat_descriptive/glossaire:term-Stylo}}}
\sphinxAtStartPar
Stylo est un éditeur de texte scientifique simplifiant la rédaction et l’édition d’articles scientifiques en Sciences Humaines et Sociales. Il est hébergé sur \sphinxurl{https://stylo.huma-num.fr}. Stylo est un logiciel libre, dont le code est disponible sur github. Il est développé developpé depuis 2020 par la Chaire de recherche du Canada sur les écritures numériques, avec le soutien de l’Université de Montréal, d’Érudit, du Centre de recherche interuniversitaire sur les humanités numériques et de l’IR* Huma\sphinxhyphen{}Num dans le cadre l’entente CRIHN\sphinxhyphen{}HN, sous licence GPL\sphinxhyphen{}3.0.

\end{description}

\sphinxstepscope


\chapter{Bibliographie}
\label{\detokenize{biblio:bibliographie}}\label{\detokenize{biblio::doc}}
\begin{sphinxthebibliography}{Gro}
\bibitem[Gro]{biblio:id3}
\sphinxAtStartPar
Maarten Grootendorst. BERTopic: Neural topic modeling with a class\sphinxhyphen{}based TF\sphinxhyphen{}IDF procedure. URL: \sphinxurl{http://arxiv.org/abs/2203.05794} (visited on 2025\sphinxhyphen{}04\sphinxhyphen{}11), \sphinxhref{https://arxiv.org/abs/2203.05794}{arXiv:2203.05794}, \sphinxhref{https://doi.org/10.48550/arXiv.2203.05794}{doi:10.48550/arXiv.2203.05794}.
\end{sphinxthebibliography}







\renewcommand{\indexname}{Index}
\printindex
\end{document}